\documentclass[10pt,a4paper]{amsart}
\usepackage[margin=19mm]{geometry}
\usepackage{amsmath,amssymb,mathtools}
\usepackage[scr=rsfs,cal=euler]{mathalfa}
\usepackage{xfrac}
\usepackage[unicode,hidelinks,bookmarks=false]{hyperref}
\newcommand{\ie}{i.e.\ }
\newcommand{\cf}{cf.\ }
\newcommand{\resp}{respectively\ }
\newcommand{\Subsection}{\subsection}
\providecommand{\nobreakdash}{\nobreak\hbox{-}}
\setlength{\emergencystretch}{2em}
\multlinegap0pt
\usepackage{bm}
\usepackage{amssymb}
\newtheorem{lemma}{Lemma}
\newcommand{\CC}{\mathbb{C}}
\newcommand{\RR}{\mathbb{R}}
\title{\sc Strong-coupling asymptotics for \\ the anisotropic Rabi model}
\author{Masao Hirokawa, Fumio Hiroshima, DongYun Lee}
\date{Source: arXiv:2510.20201v4 (CC BY 4.0)}
\begin{document}
\maketitle

\noindent\emph{Source and licence.} \sc Strong-coupling asymptotics for \\ the anisotropic Rabi model. Version arXiv:2510.20201v4 (CC BY 4.0), distributed by arXiv under the Creative Commons Attribution 4.0 International licence, http://creativecommons.org/licenses/by/4.0/, as recorded in the arXiv licence metadata for this exact version and shown on the abstract page https://arxiv.org/abs/2510.20201v4. Full licence notice: \texttt{licenses/arxiv-2510.20201-CC-BY-4.0.txt}.\par\smallskip

\noindent\textit{Editorial scope.} Counted unit: the article's lemma lem:diag\_conv (source0.tex lines 974--981). The article's complete original proof of it is delivered with it (source0.tex lines 983--1046, one \texttt{proof} environment of 64 lines). No mathematics is written, repaired or re-derived: the statement and the proof are the article's own lines, in the article's own notation, and no retained line is substituted or re-flowed.  No further source-local context is needed: every cross-reference of every retained line resolves inside the two delivered spans.  Nothing was omitted from any retained span, so the package carries no omission record.  No retained line contains a citation, so the wrapper carries no bibliography and no bibliography entry.  No retained line contains a figure environment, an image-inclusion command or drawing code, so no image is delivered and the package declares no asset dependency.  Editorial formatting only, in addition: an AMS \texttt{amsart} preamble that restates the article's own declarations and macros used by the retained lines, namely the environments \texttt{lemma} on separate counters, together with the standard packages those lines need.  The retained environments print in this document as Lemma 1 in order of appearance, whereas the article prints the same items with the numbering of its own full document.  The complete native source of the article as distributed in the arXiv e-print for this exact version is bundled unchanged as \texttt{sources/187631/source0.tex}.
\begin{lemma}\label{lem:diag_conv}
Assume $\alpha\ne0$ and $|\alpha|\ge|\beta|$. Let $z\in\CC\setminus\RR$. 
Then 
\begin{align*}
&\operatorname*{s-lim}_{|g|\to\infty} A_g(z)=(\hat{h}_\infty-z)^{-1},\\
&\operatorname*{s-lim}_{|g|\to\infty} D_g(z)=(\hat{h}_\infty-z)^{-1}.
\end{align*}
\end{lemma}

\begin{proof}
By symmetry, it suffices to prove the strong convergence for $A_g(z)$, as the proof for $D_g(z)$ is identical upon replacing $X_g$ with $X_g^\ast$.
Since $\sup_{g} \| A_g(z) \| < \infty$, ${\mathscr S}(\RR) \subset \bigcap_g D(A_g(z)^{-1}) \cap D(\hat{h}_\infty)$, and ${\mathscr S}(\RR)$ is a core for $h_{\infty} - z$, it suffices to establish the following limit on ${\mathscr S}(\RR)$:
\begin{align*}
    \lim_{|g|\to\infty} X_g (h - z)^{-1} X_g^\ast \psi 
    = \frac{\beta^2}{2\alpha^2} p^2 \psi, 
    \quad \psi \in {\mathscr S}(\RR).
\end{align*}
For any $\psi \in {\mathscr S}(\RR)$, we can write
\begin{align*}
    X_g (h - z)^{-1} X_g^\ast \psi
    = (-\triangle + ig\beta p) (h_g - z)^{-1} (-\triangle - ig\beta p) \psi,
\end{align*}
where 
\begin{align}h_g = e^{-2ig\alpha p} h e^{2ig\alpha p} = \frac12 (p^2 + (q - 2g\alpha)^2) = h - 2g\alpha q + 2g^2\alpha^2.\end{align} 
Thus our task reduces to evaluating the asymptotic behavior of
\begin{align*}
    g^n (h_g - z)^{-1}\psi, \quad g^n p (h_g - z)^{-1} \psi \quad (n = 0, 1, 2)
\end{align*}
as $|g| \to \infty$ for $\psi \in {\mathscr S}(\RR)$. Applying the second resolvent identity, we have
\begin{align*}
    \left( (h_g - z)^{-1} - \frac{1}{2g^2\alpha^2} \right) \psi
    = -(h_g - z)^{-1} (h - z - 2g\alpha q) \frac{1}{2g^2\alpha^2} \psi,
\end{align*}
which yields the identity
\begin{align}\label{n=0eq}
    (h_g - z)^{-1} \psi
    = \frac{1}{2g^2\alpha^2} \psi 
    - \frac{1}{2g^2\alpha^2} (h_g - z)^{-1} (h - z) \psi
    + \frac{1}{g\alpha} (h_g - z)^{-1} q\psi.
\end{align}
Since $\| (h_g - z)^{-1} \| = \| (h - z)^{-1} \| < \infty$, identity \eqref{n=0eq} immediately implies
\begin{align}\label{n=0lim}
    \lim_{|g|\to\infty} (h_g - z)^{-1} \psi = 0, 
    \quad 
    \psi \in {\mathscr S}(\RR).
\end{align}
Multiplying both sides of \eqref{n=0eq} by $g$ and combining with \eqref{n=0lim}, we obtain
\begin{align}\label{n=1lim}
    \lim_{|g|\to\infty} g (h_g - z)^{-1} \psi = 0, 
    \quad 
    \psi \in {\mathscr S}(\RR).
\end{align}
Furthermore, multiplying both sides of \eqref{n=0eq} by $g^2$ and using \eqref{n=0lim} and \eqref{n=1lim}, we get
\begin{align*}
    \lim_{|g|\to\infty} g^2 (h_g - z)^{-1} \psi = \frac{1}{2\alpha^2} \psi, 
    \quad 
    \psi \in {\mathscr S}(\RR).
\end{align*}
We emphasize that $g^n (h_g - z)^{-1}, n = 1, 2$ cannot converge strongly on the whole $L^2(\RR)$ since they lack uniform boundedness. We can only obtain strong convergence on a proper subspace. Next, multiplying $p$ from the left onto both sides of \eqref{n=0eq}, we have
\begin{align*}
    p (h_g - z)^{-1} \psi
    = \frac{1}{2g^2\alpha^2} p\psi 
    - \frac{1}{2g^2\alpha^2} p (h_g - z)^{-1} (h - z) \psi
    + \frac{1}{g\alpha} p (h_g - z)^{-1} q\psi.
\end{align*}
Since $\| p (h_g - z)^{-1} \| = \| p (h - z)^{-1} \| < \infty$, a similar argument yields that
\begin{align*}
    \lim_{|g|\to\infty} p (h_g - z)^{-1} \psi = \lim_{|g|\to\infty} g p (h_g - z)^{-1} \psi = 0, 
    \quad \text{and} \quad
    \lim_{|g|\to\infty} g^2 p (h_g - z)^{-1} \psi = \frac{1}{2\alpha^2} p\psi
\end{align*}
for all $\psi \in {\mathscr S}(\RR)$. Combining these asymptotic limits completes the proof for $A_g(z)$.
\end{proof}

\end{document}
